%% Copernicus Publications Manuscript Preparation Template for LaTeX Submissions
%% ---------------------------------
%% This template should be used for copernicus.cls
%% The class file and some style files are bundled in the Copernicus Latex Package, which can be downloaded from the different journal webpages.
%% For further assistance please contact Copernicus Publications at: production@copernicus.org
%% https://publications.copernicus.org/for_authors/manuscript_preparation.html


%% Please use the following documentclass and journal abbreviations for preprints and final revised papers.

%% 2-column papers and preprints
\documentclass[gmd, manuscript]{copernicus}



%% Journal abbreviations (please use the same for preprints and final revised papers)


% Advances in Geosciences (adgeo)
% Advances in Radio Science (ars)
% Advances in Science and Research (asr)
% Advances in Statistical Climatology, Meteorology and Oceanography (ascmo)
% Aerosol Research (ar)
% Annales Geophysicae (angeo)
% Archives Animal Breeding (aab)
% Atmospheric Chemistry and Physics (acp)
% Atmospheric Measurement Techniques (amt)
% Biogeosciences (bg)
% Climate of the Past (cp)
% DEUQUA Special Publications (deuquasp)
% Earth Surface Dynamics (esurf)
% Earth System Dynamics (esd)
% Earth System Science Data (essd)
% E&G Quaternary Science Journal (egqsj)
% EGUsphere (egusphere) | This is only for EGUsphere preprints submitted without relation to an EGU journal.
% Earth Observation (eo)
% European Journal of Mineralogy (ejm)
% Geochronology (gchron)
% Geographica Helvetica (gh)
% Geoscience Communication (gc)
% Geoscientific Instrumentation, Methods and Data Systems (gi)
% Geoscientific Model Development (gmd)
% History of Geo- and Space Sciences (hgss)
% Hydrology and Earth System Sciences (hess)
% Journal of Bone and Joint Infection (jbji)
% Journal of Environmentally Compatible Air Transport System (jecats)
% Journal of Micropalaeontology (jm)
% Journal of Sensors and Sensor Systems (jsss)
% Magnetic Resonance (mr)
% Mechanical Sciences (ms)
% Natural Hazards and Earth System Sciences (nhess)
% Nonlinear Processes in Geophysics (npg)
% Ocean Science (os)
% Polarforschung - Journal of the German Society for Polar Research (polf)
% Proceedings of the International Association of Hydrological Sciences (piahs)
% Proceedings of the International Ocean Drilling Programme (piodp)
% Safety of Nuclear Waste Disposal (sand)
% Scientific Drilling (sd)
% SOIL (soil)
% Solid Earth (se)
% State of the Planet (sp)
% The Cryosphere (tc)
% Weather and Climate Dynamics (wcd)
% Web Ecology (we)
% Wind Energy Science (wes)


%% \usepackage commands included in the copernicus.cls:
%\usepackage[german, english]{babel}
%\usepackage{tabularx}
%\usepackage{cancel}
%\usepackage{multirow}
%\usepackage{supertabular}
%\usepackage{algorithmic}
%\usepackage{algorithm}
%\usepackage{amsthm}
%\usepackage{float}
%\usepackage{subfig}
%\usepackage{rotating}


\usepackage{pdfpages}
\usepackage{booktabs}
\usepackage{pdflscape}

\graphicspath{{./}}
\emergencystretch=2em

\begin{document}

\title{Incompressible Stokes Flow Annulus Benchmark Using Underworld3}


% \Author[affil]{given_name}{surname}

\Author[1][gthyagi@gmail.com]{Thyagarajulu}{Gollapalli} %% correspondence author
\Author[1]{Underworld}{Team}

\affil[1]{Research School of Earth Sciences, Australian National University, Canberra, ACT, Australia}

%% The [] brackets identify the author with the corresponding affiliation. 1, 2, 3, etc. should be inserted.

%% If an author is deceased, please add \deceased[$Deceased date if applicable$]{$Author number$} (e.g. \deceased[13 November 2015]{2}) at the end of the affiliations. The author number depends on the placement of the author in the author list, e.g. the third author has number 3.


%% If authors contributed equally, please add \equalcontrib{$Author numbers$} (e.g. \equalcontrib{1,3}) at the end of the affiliations. The author number depends on the placement of the author in the author list, e.g. the third author has number 3.




\runningtitle{Annulus Stokes flow benchmark in Underworld3}

\runningauthor{Gollapalli}





\received{}
\pubdiscuss{} %% only important for two-stage journals
\revised{}
\accepted{}
\published{}

%% These dates will be inserted by Copernicus Publications during the typesetting process.


\firstpage{1}

\maketitle



\begin{abstract}
We reproduce the annulus Stokes benchmark of \citet{thieulot2018} (TP2018) using
Underworld3 (UW3). The benchmark describes incompressible isoviscous Stokes flow
in a cylindrical annulus with tangential velocity boundary conditions imposed on
both the inner and outer boundaries, driven by radially directed gravity and
analytically prescribed buoyancy forcing that generates convection cells
parameterised by the harmonic number \(k\). We investigate the accuracy and
convergence behaviour of several mixed finite-element discretisations, including
\(P_1\times P_0\), \(P_1\times P_1\), \(P_2\times P_1\), and
\(P_3\times P_2\). Numerical solutions are evaluated against the analytical
velocity and pressure fields using relative \(L_2\)-norm errors. The higher-order
Taylor--Hood discretisations reproduce the expected convergence hierarchy for
smooth Stokes solutions, with approximately third-order velocity and second-order
pressure convergence for \(P_2\times P_1\), and fourth-order velocity and
third-order pressure convergence for \(P_3\times P_2\). In contrast,
lower-order discretisations exhibit reduced robustness, particularly in pressure
recovery. These results demonstrate that UW3 accurately reproduces analytical
Stokes solutions in curved geometries and provides a robust framework for future
annulus geodynamics simulations.
\end{abstract}


% \copyrightstatement{TEXT} %% This section is optional and can be used for copyright transfers.


\introduction

Accurate numerical solution of the incompressible Stokes equations is fundamental
to modelling mantle convection, lithospheric deformation, subduction dynamics,
and planetary evolution. Many geodynamic applications require simulations in
curved geometries, such as annulus, spherical cap, and full spherical shell.
These geometries introduce additional numerical challenges associated with
curved boundaries, pressure nullspaces, coordinate-dependent forcing terms, and
tangential boundary conditions.

Reproducing analytical benchmarks remains one of the most rigorous approaches
for assessing the correctness of numerical implementations. Unlike code-to-code
comparisons, analytical benchmarks enable direct quantification of
discretisation errors and convergence rates. They are particularly valuable for
geodynamic Stokes problems because they simultaneously assess
velocity--pressure coupling, incompressibility preservation, boundary condition
implementation, and the accurate representation of curved geometries
\citep{kramer2021,davies2022}.

The annulus benchmark introduced by TP2018 provides an exact analytical solution
for incompressible isoviscous Stokes flow in a cylindrical shell subjected to
radially directed gravity. The benchmark was specifically designed for
geodynamics applications and has since been widely used for verification of
geodynamic codes within the community \citep{gassmoller2024}. Unlike classical
Cartesian manufactured solutions, the benchmark incorporates curved inner and
outer boundaries, non-trivial pressure distributions along the boundaries,
coordinate-dependent buoyancy forcing, tangential velocity boundary conditions,
and controllable azimuthal complexity through the harmonic parameter \(k\).
These characteristics make the benchmark particularly well suited for assessing
finite-element formulations in annulus and curved-domain mantle convection
problems.

UW3 is a next-generation geodynamics framework built around PETSc and symbolic
problem specification in Python. In this study, we reproduce the TP2018 annulus
benchmark to verify the curved-domain finite-element implementations in UW3. We
evaluate the convergence and accuracy of several mixed finite-element
discretisations, including both lower-order and higher-order
velocity--pressure element pairs. Numerical results are compared against the
analytical velocity and pressure fields using \(L_2\)-norm errors and
boundary pressure diagnostics. Particular attention is given to the convergence
behaviour of Taylor--Hood discretisations and the performance of lower-order and
equal-order element formulations.

\section{Governing Equations}

We consider steady state incompressible Stokes equations:
\begin{subequations}
\begin{align}
  -\nabla \cdot \boldsymbol{\sigma} &= \rho \mathbf{g}, \\
  \nabla \cdot \mathbf{u} &= 0.
\end{align}
\end{subequations}
The stress tensor is
\begin{equation}
  \boldsymbol{\sigma} = -p\mathbf{I} + 2\eta\boldsymbol{\varepsilon}(\mathbf{u}),
\end{equation}
with viscosity \(\eta\) and strain-rate tensor
\begin{equation}
  \boldsymbol{\varepsilon}(\mathbf{u})
  =
  \frac{1}{2}\left(\nabla \mathbf{u} + (\nabla \mathbf{u})^T\right).
\end{equation}
where \(\mathbf{u}\) is velocity, \(p\) is pressure, \(\rho\) is density, and \(\mathbf{g}\) is gravitational acceleration.

\section{Analytical Benchmark Solution}

Following \citet{thieulot2018}, the analytical velocity field in polar coordinates is defined as
\begin{subequations}
\begin{align}
  v_\theta(r,\theta) &= f(r)\cos(k\theta), \\
  v_r(r,\theta) &= k g(r)\sin(k\theta).
\end{align}
\end{subequations}
where
\begin{subequations}
\begin{align}
  f(r) &= Ar+\frac{B}{r}, \\
  g(r) &= \frac{A}{2}r+\frac{B}{r}\ln r+\frac{C}{r}.
\end{align}
\end{subequations}
The constants \(A\) and \(B\) are chosen such that radial velocity vanishes at both annulus boundaries:
\begin{subequations}
\begin{align}
  A &= -C\frac{2(\ln R_1-\ln R_2)}{R_2^2\ln R_1-R_1^2\ln R_2}, \\
  B &= -C\frac{R_2^2-R_1^2}{R_2^2\ln R_1-R_1^2\ln R_2}.
\end{align}
\end{subequations}
The pressure field is
\begin{equation}
  p(r,\theta)=k h(r)\sin(k\theta)+\rho_0(R_2-r),
\end{equation}
where
\begin{equation}
  h(r)=\frac{2g(r)-f(r)}{r}.
\end{equation}
The corresponding density field is
\begin{equation}
  \rho(r,\theta)=kM(r)\sin(k\theta)+\rho_0,
\end{equation}
with
\begin{equation}
  M(r)=g''(r)-\frac{g'(r)}{r}
  -\frac{g(r)}{r^2}(k^2-1)
  +\frac{f(r)}{r^2}
  +\frac{f'(r)}{r}.
\end{equation}
The benchmark calculations use the standard parameter values
\[
  R_1=1,
  \qquad
  R_2=2,
  \qquad
  C=-1,
  \qquad
  \rho_0=0.
\]
where \(R_1\) is the inner radius and \(R_2\) is the outer radius.
Gravity acts radially inward:
\begin{equation}
  \mathbf{g} = -\mathbf{e}_r.
\end{equation}
The benchmark assumes constant viscosity \(\eta=1\) and prescribes analytical velocity on boundaries, and density field in the domain.
Increasing the harmonic parameter \(k\) increases azimuthal complexity and produces \(2k\) convection cells in the annulus.
This enables systematic evaluation of numerical performance across increasingly oscillatory flow structures.

\begin{figure}[H]
    \centering
    \includegraphics[width=0.98\linewidth]{figure_1_thieulot_annulus_analytical.pdf}
    \caption{Analytical density, velocity, and pressure fields. Columns correspond to increasing \(k\), which increases the number of convection cells in the annulus.}
    \label{fig:thieulot_annulus_fields}
\end{figure}

\section{Numerical Methodology}

\subsection{UW3 Formulation}

The benchmark is implemented in UW3 using PETSc-based finite-element Stokes system.
The governing equations are assembled symbolically using SymPy and solved through PETSc Krylov subspace methods.
A constant-pressure nullspace is supplied to PETSc to remove the pressure gauge ambiguity.
Following solution, the domain-average numerical pressure is removed prior to error evaluation.
The annulus meshes are generated using unstructured triangular elements with progressively refined nominal cell sizes.
The mesh geometry degree is set to linear for all reported calculations.

\subsection{Finite-Element Discretisations}

The following mixed finite-element pairs are evaluated:
\begin{table}[H]
\centering
\begin{tabular}{lll}
\toprule
Velocity space & Pressure space & Description \\
\midrule
\(P_1\) & \(P_0\) & Linear velocity, discontinuous constant pressure \\
\(P_1\) & \(P_1\) & Equal-order continuous discretisation \\
\(P_2\) & \(P_1\) & Taylor--Hood discretisation \\
\(P_3\) & \(P_2\) & Higher-order Taylor--Hood discretisation \\
\bottomrule
\end{tabular}
\end{table}

The quadrature degree is chosen according to the highest polynomial order present in the discretisation.
For smooth incompressible Stokes solutions, stable Taylor--Hood-type discretisations are expected to exhibit velocity convergence of order \(p+1\) and pressure convergence of order \(p\) for a \(P_{p+1}\times P_p\) discretisation \citep{girault1986,boffi2013}.

\subsection{Error Metrics}

Relative \(L_2\)-norm errors are computed as
\begin{equation}
  \|e_q\|_{L_2,\mathrm{rel}}
  =
  \left(
  \frac{\int_\Omega |q_h-q|^2\,d\Omega}
       {\int_\Omega |q|^2\,d\Omega}
  \right)^{1/2},
\end{equation}
where \(q_h\) is the numerical solution and \(q\) is the analytical solution.
Convergence rates are computed from successive mesh refinements.

add convergence rate formula. add absolute l2 norms.

\section{Results}

\subsection{Convergence Behaviour}

The convergence behaviour for all tested discretisations is shown in Fig.~\ref{fig:thieulot_annulus_convergence}.
The \(P_2\times P_1\) Taylor--Hood discretisation exhibits approximately third-order velocity convergence and second-order pressure convergence, consistent with theoretical expectations for smooth Stokes solutions.

The higher-order \(P_3\times P_2\) discretisation achieves approximately fourth-order velocity convergence and third-order pressure convergence on coarse-to-intermediate mesh resolutions.
At the finest resolutions, convergence rates begin to plateau, likely reflecting accumulated effects from geometry approximation, quadrature errors, and solver tolerances.

The lower-order discretisations show reduced robustness.
The \(P_1\times P_0\) pair fails to produce reliable convergence behaviour for this benchmark, while the equal-order \(P_1\times P_1\) discretisation produces acceptable velocity convergence but weaker and less consistent pressure convergence.
These observations are consistent with established finite-element theory for incompressible Stokes flow \citep{boffi2013,thieulot2022}.

\begin{figure}[H]
    \centering
    \includegraphics[width=0.98\linewidth]{figure_5_thieulot_annulus_convergence.pdf}
    \caption{Convergence of relative volume \(L_2\)-norm velocity and pressure errors for the annulus Thieulot benchmark. Reference lines indicate the expected smooth-solution rates for the higher-order element pairs.}
    \label{fig:thieulot_annulus_convergence}
\end{figure}

% Native LaTeX table content from table_figure_5_thieulot_annulus_convergence.tex.
\par\nolinenumbers
\begin{landscape}
\thispagestyle{empty}
\begin{center}
\centering
\captionof{table}{Errors and Convergence Rates for $P_1\times P_0$.}
\label{tab:annulus_thieulot_v1_p0_disc}
\small
\begin{tabular}{
r rr rr rr rr rr rr
}
\toprule
& \multicolumn{4}{c}{$k=1$} & \multicolumn{4}{c}{$k=4$} & \multicolumn{4}{c}{$k=8$} \\
\cmidrule(lr){2-5}\cmidrule(lr){6-9}\cmidrule(lr){10-13}
$h$ & $|e_v|_2$ & rate & $|e_p|_2$ & rate & $|e_v|_2$ & rate & $|e_p|_2$ & rate & $|e_v|_2$ & rate & $|e_p|_2$ & rate \\
\midrule
$1/8$ & 6.398e-02 & -- & 8.040e-01 & -- & 1.763e-01 & -- & 7.800e-01 & -- & 1.298e-01 & -- & 8.163e-01 & -- \\
$1/16$ & 6.427e-02 & -0.01 & 7.384e-01 & 0.12 & 1.775e-01 & -0.01 & 7.295e-01 & 0.10 & 1.401e-01 & -0.11 & 7.589e-01 & 0.11 \\
$1/32$ & 6.426e-02 & 0.00 & 7.220e-01 & 0.03 & 1.780e-01 & -0.00 & 7.161e-01 & 0.03 & 1.435e-01 & -0.04 & 7.403e-01 & 0.04 \\
$1/64$ & 6.426e-02 & 0.00 & 7.177e-01 & 0.01 & 1.781e-01 & -0.00 & 7.124e-01 & 0.01 & 1.445e-01 & -0.01 & 7.349e-01 & 0.01 \\
$1/128$ & 6.425e-02 & 0.00 & 7.167e-01 & 0.00 & 1.782e-01 & -0.00 & 7.114e-01 & 0.00 & 1.447e-01 & -0.00 & 7.334e-01 & 0.00 \\
$1/256$ & 6.425e-02 & 0.00 & 7.164e-01 & 0.00 & 1.782e-01 & -0.00 & 7.112e-01 & 0.00 & 1.448e-01 & -0.00 & 7.330e-01 & 0.00 \\
$1/512$ & 6.425e-02 & 0.00 & 7.163e-01 & 0.00 & 1.782e-01 & -0.00 & 7.111e-01 & 0.00 & 1.448e-01 & -0.00 & 7.329e-01 & 0.00 \\
\bottomrule
\end{tabular}
\end{center}

\begin{center}
\centering
\captionof{table}{Errors and Convergence Rates for $P_1\times P_1$.}
\label{tab:annulus_thieulot_v1_p1_cont}
\small
\begin{tabular}{
r rr rr rr rr rr rr
}
\toprule
& \multicolumn{4}{c}{$k=1$} & \multicolumn{4}{c}{$k=4$} & \multicolumn{4}{c}{$k=8$} \\
\cmidrule(lr){2-5}\cmidrule(lr){6-9}\cmidrule(lr){10-13}
$h$ & $|e_v|_2$ & rate & $|e_p|_2$ & rate & $|e_v|_2$ & rate & $|e_p|_2$ & rate & $|e_v|_2$ & rate & $|e_p|_2$ & rate \\
\midrule
$1/8$ & 1.537e-02 & -- & 4.325e+00 & -- & 2.994e-02 & -- & 3.202e+00 & -- & 3.629e-02 & -- & 3.184e+00 & -- \\
$1/16$ & 3.459e-03 & 2.15 & 1.628e+00 & 1.41 & 6.385e-03 & 2.23 & 1.625e+00 & 0.98 & 7.925e-03 & 2.19 & 1.417e+00 & 1.17 \\
$1/32$ & 8.358e-04 & 2.05 & 7.174e-01 & 1.18 & 1.293e-03 & 2.30 & 6.195e-01 & 1.39 & 1.937e-03 & 2.03 & 7.144e-01 & 0.99 \\
$1/64$ & 2.056e-04 & 2.02 & 3.115e-01 & 1.20 & 2.776e-04 & 2.22 & 2.855e-01 & 1.12 & 4.931e-04 & 1.97 & 3.210e-01 & 1.15 \\
$1/128$ & 5.095e-05 & 2.01 & 1.582e-01 & 0.98 & 6.252e-05 & 2.15 & 1.083e-01 & 1.40 & 1.243e-04 & 1.99 & 1.260e-01 & 1.35 \\
$1/256$ & 1.266e-05 & 2.01 & 5.740e-02 & 1.46 & 1.487e-05 & 2.07 & 5.231e-02 & 1.05 & 3.158e-05 & 1.98 & 6.476e-02 & 0.96 \\
\bottomrule
\end{tabular}
\end{center}

\vfill
\begin{center}\thepage\end{center}
\clearpage
\end{landscape}
\begin{landscape}
\thispagestyle{empty}

\begin{center}
\centering
\captionof{table}{Errors and Convergence Rates for $P_2\times P_1$.}
\label{tab:annulus_thieulot_v2_p1_cont}
\small
\begin{tabular}{
r rr rr rr rr rr rr
}
\toprule
& \multicolumn{4}{c}{$k=1$} & \multicolumn{4}{c}{$k=4$} & \multicolumn{4}{c}{$k=8$} \\
\cmidrule(lr){2-5}\cmidrule(lr){6-9}\cmidrule(lr){10-13}
$h$ & $|e_v|_2$ & rate & $|e_p|_2$ & rate & $|e_v|_2$ & rate & $|e_p|_2$ & rate & $|e_v|_2$ & rate & $|e_p|_2$ & rate \\
\midrule
$1/8$ & 2.218e-04 & -- & 1.087e-02 & -- & 9.781e-04 & -- & 1.944e-02 & -- & 2.731e-03 & -- & 4.174e-02 & -- \\
$1/16$ & 2.675e-05 & 3.05 & 2.530e-03 & 2.10 & 1.151e-04 & 3.09 & 4.611e-03 & 2.08 & 3.225e-04 & 3.08 & 9.354e-03 & 2.16 \\
$1/32$ & 3.388e-06 & 2.98 & 6.427e-04 & 1.98 & 1.437e-05 & 3.00 & 1.159e-03 & 1.99 & 4.016e-05 & 3.01 & 2.310e-03 & 2.02 \\
$1/64$ & 4.205e-07 & 3.01 & 1.601e-04 & 2.01 & 1.763e-06 & 3.03 & 2.869e-04 & 2.01 & 4.850e-06 & 3.05 & 5.666e-04 & 2.03 \\
$1/128$ & 5.254e-08 & 3.00 & 3.992e-05 & 2.00 & 2.190e-07 & 3.01 & 7.152e-05 & 2.00 & 6.001e-07 & 3.01 & 1.412e-04 & 2.00 \\
$1/256$ & 6.575e-09 & 3.00 & 9.979e-06 & 2.00 & 2.730e-08 & 3.00 & 1.787e-05 & 2.00 & 7.449e-08 & 3.01 & 3.524e-05 & 2.00 \\
$1/512$ & 8.632e-10 & 2.93 & 2.494e-06 & 2.00 & 3.412e-09 & 3.00 & 4.466e-06 & 2.00 & 9.287e-09 & 3.00 & 8.802e-06 & 2.00 \\
\bottomrule
\end{tabular}
\end{center}

\begin{center}
\centering
\captionof{table}{Errors and Convergence Rates for $P_3\times P_2$.}
\label{tab:annulus_thieulot_v3_p2_cont}
\small
\begin{tabular}{
r rr rr rr rr rr rr
}
\toprule
& \multicolumn{4}{c}{$k=1$} & \multicolumn{4}{c}{$k=4$} & \multicolumn{4}{c}{$k=8$} \\
\cmidrule(lr){2-5}\cmidrule(lr){6-9}\cmidrule(lr){10-13}
$h$ & $|e_v|_2$ & rate & $|e_p|_2$ & rate & $|e_v|_2$ & rate & $|e_p|_2$ & rate & $|e_v|_2$ & rate & $|e_p|_2$ & rate \\
\midrule
$1/8$ & 1.102e-05 & -- & 1.210e-03 & -- & 6.233e-05 & -- & 2.498e-03 & -- & 2.198e-04 & -- & 6.495e-03 & -- \\
$1/16$ & 6.792e-07 & 4.02 & 1.220e-04 & 3.31 & 3.811e-06 & 4.03 & 2.905e-04 & 3.10 & 1.363e-05 & 4.01 & 8.249e-04 & 2.98 \\
$1/32$ & 4.360e-08 & 3.96 & 1.392e-05 & 3.13 & 2.441e-07 & 3.96 & 3.358e-05 & 3.11 & 8.723e-07 & 3.97 & 9.664e-05 & 3.09 \\
$1/64$ & 2.729e-09 & 4.00 & 1.585e-06 & 3.13 & 1.510e-08 & 4.01 & 3.870e-06 & 3.12 & 5.323e-08 & 4.03 & 1.110e-05 & 3.12 \\
$1/128$ & 4.043e-10 & 2.75 & 1.877e-07 & 3.08 & 9.826e-10 & 3.94 & 4.601e-07 & 3.07 & 3.344e-09 & 3.99 & 1.322e-06 & 3.07 \\
\bottomrule
\end{tabular}
\end{center}

\vfill
\begin{center}\thepage\end{center}
\end{landscape}
\par\linenumbers

\subsection{Boundary Pressure Recovery}

Pressure profiles extracted along the inner and outer annulus boundaries are compared against analytical solutions in Fig.~\ref{fig:thieulot_annulus_boundary_pressure}.
The higher-order Taylor--Hood discretisations reproduce the analytical pressure distributions accurately across the annulus boundaries, including extrema near regions of high azimuthal variation.
Pressure discrepancies are most pronounced for lower-order discretisations and become increasingly visible for larger \(k\), where sharper angular pressure gradients must be resolved.

Because the benchmark includes non-trivial pressure variation directly on curved boundaries, these comparisons provide a stringent test of both pressure interpolation quality and curved-domain numerical accuracy.

\begin{figure}[H]
    \centering
    \includegraphics[width=0.92\linewidth]{figure_2_thieulot_annulus_pressure_boundary.pdf}
    \caption{Pressure along the inner and outer annulus boundaries for \(k=1,2,4\), using the \(P_2\times P_1\), \(h=1/128\) calculations. Solid curves denote the analytical solution and dashed curves denote UW3 values sampled from the numerical pressure field.}
    \label{fig:thieulot_annulus_boundary_pressure}
\end{figure}

\subsection{Implications for Spherical Geodynamics}

Although the benchmark is formulated in a two-dimensional annulus, its numerical characteristics are directly relevant to spherical-shell mantle convection calculations.
In particular, the benchmark tests curved-boundary finite-element behaviour, radial body-force implementation, pressure nullspace treatment, incompressibility preservation, and higher-order mixed discretisations.
These are all essential ingredients in regional and global spherical geodynamics models.

The strong convergence behaviour observed for higher-order Taylor--Hood discretisations therefore provides confidence in Underworld3 implementations intended for annulus and spherical-shell mantle convection applications.

\section{Discussion}

The present results demonstrate that Underworld3 reproduces the analytical annulus benchmark solution with convergence behaviour consistent with finite-element theory.
The benchmark additionally highlights several important numerical considerations.

First, higher-order mixed discretisations provide substantially improved pressure accuracy relative to lower-order formulations.
This is particularly important in geodynamics applications where pressure gradients strongly influence buoyancy-driven flow and dynamic topography.

Second, curved geometries remain sensitive to geometric approximation order.
Although linear geometry representation was sufficient to recover theoretical convergence trends for the present benchmark, higher-order curved geometry representations may become increasingly important for extremely fine meshes or strongly curved spherical-shell problems.

Third, equal-order discretisations continue to exhibit weaker pressure robustness relative to stable Taylor--Hood formulations.
This behaviour aligns with expectations from the Ladyzhenskaya--Babuska--Brezzi stability condition.

The benchmark therefore provides a valuable verification framework not only for Underworld3, but also for future development of curved-element geometries, advanced boundary-condition implementations, and spherical-shell finite-element formulations.

\conclusions

We have presented a verification study of the Thieulot--Puckett annulus Stokes benchmark using Underworld3.
The benchmark provides a stringent test of incompressible Stokes flow in curved geometries through the combination of radial gravity, coordinate-dependent forcing, tangential boundary velocity conditions, and non-trivial pressure variation on curved annulus boundaries.

The principal findings are:
\begin{enumerate}
  \item Higher-order Taylor--Hood discretisations reproduce the expected smooth-solution convergence hierarchy.
  \item The \(P_2\times P_1\) discretisation achieves approximately third-order velocity and second-order pressure convergence.
  \item The \(P_3\times P_2\) discretisation achieves approximately fourth-order velocity and third-order pressure convergence before fine-resolution error saturation.
  \item Lower-order discretisations exhibit weaker pressure robustness and reduced convergence quality.
  \item Underworld3 accurately reproduces analytical Stokes solutions in curved annular geometries.
\end{enumerate}

These results establish the Thieulot--Puckett annulus benchmark as an effective verification framework for Underworld3 annulus and spherical-shell geodynamics applications.

%% The following commands are for the statements about the availability of data sets and/or software code corresponding to the manuscript.
%% It is strongly recommended to make use of these sections in case data sets and/or software code have been part of your research the article is based on.

% \codeavailability{TEXT} %% use this section when having only software code available


% \dataavailability{TEXT} %% use this section when having only data sets available


% \codedataavailability{TEXT} %% use this section when having data sets and software code available


% \sampleavailability{TEXT} %% use this section when having geoscientific samples available


% \videosupplement{TEXT} %% use this section when having video supplements available


% \appendix
% \section{}    %% Appendix A

% \subsection{}     %% Appendix A1, A2, etc.


% \noappendix       %% use this to mark the end of the appendix section. Otherwise the figures might be numbered incorrectly (e.g. 10 instead of 1).

%% Regarding figures and tables in appendices, the following two options are possible depending on your general handling of figures and tables in the manuscript environment:

%% Option 1: If you sorted all figures and tables into the sections of the text, please also sort the appendix figures and appendix tables into the respective appendix sections.
%% They will be correctly named automatically.

%% Option 2: If you put all figures after the reference list, please insert appendix tables and figures after the normal tables and figures.
%% To rename them correctly to A1, A2, etc., please add the following commands in front of them:

% % \appendixfigures  %% needs to be added in front of appendix figures

% % \appendixtables   %% needs to be added in front of appendix tables

%% Please add \clearpage between each table and/or figure. Further guidelines on figures and tables can be found below.



% \authorcontribution{TEXT} %% this section is mandatory

% \competinginterests{TEXT} %% this section is mandatory even if you declare that no competing interests are present

% \disclaimer{TEXT} %% optional section

% \begin{acknowledgements}
% TEXT
% \end{acknowledgements}




%% REFERENCES

%% The reference list is compiled as follows:

\bibliographystyle{copernicus}
\bibliography{thieulot_annulus_benchmark_report_refs}

%% Since the Copernicus LaTeX package includes the BibTeX style file copernicus.bst,
%% authors experienced with BibTeX only have to include the following two lines:
%%
%% \bibliographystyle{copernicus}
%% \bibliography{example.bib}
%%
%% URLs and DOIs can be entered in your BibTeX file as:
%%
%% URL = {http://www.xyz.org/~jones/idx_g.htm}
%% DOI = {10.5194/xyz}


%% LITERATURE CITATIONS
%%
%% command                        & example result
%% \citet{jones90}|               & Jones et al. (1990)
%% \citep{jones90}|               & (Jones et al., 1990)
%% \citep{jones90,jones93}|       & (Jones et al., 1990, 1993)
%% \citep[p.~32]{jones90}|        & (Jones et al., 1990, p.~32)
%% \citep[e.g.,][]{jones90}|      & (e.g., Jones et al., 1990)
%% \citep[e.g.,][p.~32]{jones90}| & (e.g., Jones et al., 1990, p.~32)
%% \citeauthor{jones90}|          & Jones et al.
%% \citeyear{jones90}|            & 1990



%% FIGURES

%% When figures and tables are placed at the end of the MS (article in one-column style), please add \clearpage
%% between bibliography and first table and/or figure as well as between each table and/or figure.

% The figure files should be labelled correctly with Arabic numerals (e.g. fig01.jpg, fig02.png).


%% ONE-COLUMN FIGURES

%%f
%\begin{figure}[t]
%\includegraphics[width=8.3cm]{FILE NAME}
%\caption{TEXT}
%\end{figure}
%
%%% TWO-COLUMN FIGURES
%
%%f
%\begin{figure*}[t]
%\includegraphics[width=12cm]{FILE NAME}
%\caption{TEXT}
%\end{figure*}
%
%
%%% TABLES
%%%
%%% The different columns must be seperated with a & command and should
%%% end with \\ to identify the column brake.
%
%%% ONE-COLUMN TABLE
%
%%t
%\begin{table}[t]
%\caption{TEXT}
%\begin{tabular}{column = lcr}
%\tophline
%
%\middlehline
%
%\bottomhline
%\end{tabular}
%\belowtable{} % Table Footnotes
%\end{table}
%
%%% TWO-COLUMN TABLE
%
%%t
%\begin{table*}[t]
%\caption{TEXT}
%\begin{tabular}{column = lcr}
%\tophline
%
%\middlehline
%
%\bottomhline
%\end{tabular}
%\belowtable{} % Table Footnotes
%\end{sidewaystable*}
%
%%% LANDSCAPE TABLE
%
%%t
%\begin{sidewaystable*}[t]
%\caption{TEXT}
%\begin{tabular}{column = lcr}
%\tophline
%
%\middlehline
%
%\bottomhline
%\end{tabular}
%\belowtable{} % Table Footnotes
%\end{table*}
%
%
%%% MATHEMATICAL EXPRESSIONS
%
%%% All papers typeset by Copernicus Publications follow the math typesetting regulations
%%% given by the IUPAC Green Book (IUPAC: Quantities, Units and Symbols in Physical Chemistry,
%%% 2nd Edn., Blackwell Science, available at: http://old.iupac.org/publications/books/gbook/green_book_2ed.pdf, 1993).
%%%
%%% Physical quantities/variables are typeset in italic font (t for time, T for Temperature)
%%% Indices which are not defined are typeset in italic font (x, y, z, a, b, c)
%%% Items/objects which are defined are typeset in roman font (Car A, Car B)
%%% Descriptions/specifications which are defined by itself are typeset in roman font (abs, rel, ref, tot, net, ice)
%%% Abbreviations from 2 letters are typeset in roman font (RH, LAI)
%%% Vectors are identified in bold italic font using \vec{x}
%%% Matrices are identified in bold roman font
%%% Multiplication signs are typeset using the LaTeX commands \times (for vector products, grids, and exponential notations) or \cdot
%%% The character * should not be applied as mutliplication sign
%
%
%%% EQUATIONS
%
%%% Single-row equation
%
%\begin{equation}
%
%\end{equation}
%
%%% Multiline equation
%
%\begin{align}
%& 3 + 5 = 8\\
%& 3 + 5 = 8\\
%& 3 + 5 = 8
%\end{align}
%
%
%%% MATRICES
%
%\begin{matrix}
%x & y & z\\
%x & y & z\\
%x & y & z\\
%\end{matrix}
%
%
%%% ALGORITHM
%
%\begin{algorithm}
%\caption{...}
%\label{a1}
%\begin{algorithmic}
%...
%\end{algorithmic}
%\end{algorithm}
%
%
%%% CHEMICAL FORMULAS AND REACTIONS
%
%%% For formulas embedded in the text, please use \chem{}
%
%%% The reaction environment creates labels including the letter R, i.e. (R1), (R2), etc.
%
%\begin{reaction}
%%% \rightarrow should be used for normal (one-way) chemical reactions
%%% \rightleftharpoons should be used for equilibria
%%% \leftrightarrow should be used for resonance structures
%\end{reaction}
%
%
%%% PHYSICAL UNITS
%%%
%%% Please use \unit{} and apply the exponential notation (e.g. 20\,\unit{W\,m^{-2}})


\end{document}
